% SPDX-License-Identifier: LPPL-1.3c
% Copyright (C) 2026 Christophe Jorssen
% Author and Current Maintainer: Christophe Jorssen <christophe.jorssen@gmail.com>
% LPPL maintenance status: maintained
% This file is part of luafemm. See LICENSE and LICENSES.md for its terms.
\section{FEMM problem and solution interchange}
\label{sec:interchange}
The interchange described here applies to |field model=open|. Ideal-mode
floating-window constraints cannot be represented by this subset; FEM/ANS
export and FEM import in ideal mode report an error.

A drawing can now travel between your TeX document and a magnetic FEMM problem.
There are three distinct operations: read a preprocessor |.fem| file, write a
problem for another solver, and write the solution already computed by Lua as
an |.ans| file. None requires FEMM, xfemm or a native library at document runtime.
The optional xfemm tools used to validate this implementation are developer tools.
This chapter describes the supported planar, real, static subset of magnetic
format~4.0. It does not promise that every physical feature of FEMM is implemented.

\subsection{First, keep drawing curves}

Imagine that one side of a conductor is a Bezier curve and another conductor
is a rotated ellipse. We need not replace these paths by hand-written circular
arcs. PGF evaluates the drawing; luafemm approximates its curves by short straight
segments before generating finite elements. This already happens in ordinary
luafemm calculations. The FEM writer uses those computational interfaces.

There are therefore three geometric representations to keep in mind:
\begin{enumerate}
\item The evaluated source path contains the PGF lines and cubic curves, or
      the segments and circular arcs read from a FEM file. Native captured
      soft-path data and their coordinate frame are retained in the model.
\item The computational geometry is polygonal. Its faces carry materials and
      excitations, and its edges carry boundary conditions. It is independent
      of the interior triangle spacing.
\item The solution consists of the actual triangles and nodal potential $A_z$.
      ANS export preserves these data without remeshing or projection.
\end{enumerate}

The source TeX program, node anchors, text, colours and macros are not encoded
in a FEM file. A round trip preserves the represented physical problem, not the
original drawing program. Similarly, an imported circular arc may be preserved
by the passive Lua document writer, whereas the computational exporter writes
its polygonal approximation. No circular-arc fitting is performed on arbitrary
TikZ curves.

Start with a modest |curve tolerance| and |mesh size|. The first controls curve
approximation, the second triangle spacing. Reducing the triangle size cannot
recover curvature already removed by polygonalisation. As the drawing approaches
its final form, reduce each independently and compare the field on the contour
of interest. A narrow curved air gap deserves particular attention: its width
must be resolved geometrically before refining triangles inside it.
The chord tolerance is not a global error estimate for the magnetic field;
PGF coordinate precision, snapping and constructed intersections also matter.

Here is a complete example. The dotted interfaces expose the polygonal geometry;
the last drawing comes from importing the exported FEM problem and solving again.
The graph still refers to the original named profile.
\luafemmexample{interchange-curves}

\subsection{Reading a preprocessor problem}

\begin{key}{/tikz/femm/import=\marg{options}}
Create a model from a magnetic FEM file at the beginning of the picture, after
picture options have been collected. The option family is |/femm/import|.
This key installs the same millimetre coordinate convention as |femm/problem|.
It can be combined with |femm/problem| to override numerical settings explicitly.
Importing does not itself draw, mesh or solve anything.
\end{key}
\begin{key}{/femm/import/file=\meta{filename} (initially empty)}
Required input filename, including |.fem|. The bridge first uses TeX's normal
file search, then checks the output directory for a file generated earlier in
the same document, and finally uses the supplied path. The passive Lua API uses
ordinary working-directory paths. Imported content is never executed as Lua or TeX.
\end{key}

Only explicitly set problem keys override imported settings. In particular,
unspecified luafemm defaults do not replace the file's depth or solver precision.
The domain bounds come from the imported geometry; explicitly supplying |xmin|,
|xmax|, |ymin| or |ymax| together with import is an error. Source length units are
converted directly into model millimetres without a TikZ drawing/recapture step.
The Lua API can instead request another model unit.

The initial imported mesh size is one twelfth of the largest bounding-box extent.
A block-label mesh control supplies a local spacing hint in converted model units.
FEMM's diameter-like area control, minimum-angle setting and smart-mesh heuristics
are retained in passive metadata but do not imply identical meshing algorithms.
Use the documented luafemm mesh and quality keys explicitly for the Lua solve.

An imported model's physical topology is immutable in this version. Add ordinary
TikZ annotations, measurement contours and plots freely, but edit physical
boundaries/material regions in the FEM source before reimporting. Additional
|femm/region| declarations and rectangular side-condition overrides are rejected
rather than being silently ignored. Native TikZ models retain their existing
paths, components, painting precedence and fill-rule interface.

\subsection{Showing the computational geometry}

\begin{key}{/tikz/pics/femm geometry}
Draw the compiled physical interfaces of the current model. This picture does
not register another material region and does not start a magnetic solve.
For a native TikZ model it displays the polygonal interfaces used by the
computational exporter, including the exterior rectangle. For an imported model
it includes the boundaries of true exclusions.
\end{key}
\begin{stylekey}{/tikz/femm geometry}
Base drawing style, initially |draw=black!65,line width=.3pt|.
\end{stylekey}
\begin{key}{/tikz/femm/geometry style=\marg{TikZ options}}
Set |femm current geometry|, applied after the base style. Use this to inspect
interfaces in another colour or with a dotted line without changing the physics.
\end{key}
\begin{stylekey}{/tikz/femm current geometry}
Additional geometry style, initially empty.
\end{stylekey}
\begin{stylekey}{/tikz/femm group/\meta{integer}}
Optional style for an imported segment group. It is tried after the base and
current geometry styles. For example, |femm group/3/.style={red,thick}| highlights
group~3. Undeclared group styles have no effect. Imported names are never turned
into executable TeX control sequences.
\end{stylekey}

\subsection{Writing a problem or its existing solution}

\begin{key}{/tikz/femm/export=\marg{options}}
Append an export request to the picture's end hook. Repeated keys produce repeated
exports, bound to that picture's model. All physical declarations precede export.
A FEM request compiles geometry but does not mesh or solve. An ANS request requires
a converged solution at that point: use a field picture or |\femmsolve| earlier.
A missing solution is an error, not an implicit request for an expensive solve.
\end{key}
\begin{key}{/femm/export/format=\meta{choice} (initially fem)}
A PGF |.is choice| selector with choices |fem| and |ans|. FEM writes a magnetic
problem, ANS writes that problem followed by the existing solution.
\end{key}
\begin{key}{/femm/export/file=\meta{filename} (initially empty)}
Required destination, including its extension. Relative paths follow TeX's
|--output-directory|, as cache files do. Parent directories must already exist.
The writer validates and serializes first, writes a temporary file beside the
destination, then replaces the destination after success. Use a different filename
from the imported source. Repeated exports to the same destination replace it.
\end{key}
\begin{key}{/femm/export/model=\meta{name} (initially empty)}
Select the model for the explicit export command. Empty means the current Lua
model. Picture-level export supplies its bound model automatically, so the option
is intended for delayed command use. Named models remain available after their
picture has ended.
\end{key}
\begin{command}{\femmexport\opt{\oarg{options}}}
Immediately export using the same |/femm/export| family. This is useful after a
profile outside the source picture has triggered its solve. Use |model| explicitly
when subsequent pictures have created other models. It never starts a solve.
\end{command}
\begin{key}{/femm/problem/name=\meta{identifier} (initially empty)}
Register the picture model under a unique name for delayed export. Duplicate
names are errors. Prefer descriptive, nonnumeric names; internal picture IDs
are reserved. Naming does not change the numerical problem or cache identity.
\end{key}
\begin{key}{/femm/problem/depth=\meta{positive number} (initially unset)}
Planar thickness in model millimetres. A native model must specify it before
FEM or ANS export. Imported depth is converted from the declared file unit.
Depth does not change the planar $A_z$ solve, but it matters to FEMM's integrated
quantities. It is included in the solution export signature.
\end{key}
\begin{key}{/femm/problem/interpolation=\meta{choice} (initially linear)}
Choose |linear|, the existing piecewise-linear $H(B)$ law with vacuum-slope
extrapolation, or |femm|, a natural cubic law with FEMM-compatible smoothing and
endpoint-slope extrapolation. Imported nonlinear materials default to |femm|
unless this key explicitly overrides it. Set it before declaring materials.
Linear materials are unaffected.
\end{key}

Nonlinear ANS export requires |interpolation=femm| for every nonlinear material
used by the mesh. The postprocessor reconstructs $B$ from nodal $A$, then computes
$H$ from its own material law; preserving the B--H sample table alone is insufficient.
The writer therefore rejects a nonlinear solution computed with the old law.
FEM export remains available with either policy, but a subsequent FEMM solve
uses FEMM's constitutive law and can produce a different solution.

The static ANS writer includes zero-based element-to-label references, one current
record per block label and explicit empty periodic/air-gap trailers. Prescribed
current density is encoded in effective material records. The output uses the
non-incremental four-column element dialect read by the audited postprocessors.
It contains the actual converged mesh and physical potential in tesla metres.
It is not a replacement for the private |.lfc| cache and cannot currently be
imported as a solved luafemm model.

\subsection{The order of operations}

\begin{enumerate}
\item Collect picture options and create or import exactly one model. Apply only
      explicit overrides to imported numerical settings.
\item Declare the native materials and physical paths, or draw a view of imported
      geometry. Declare native boundary conditions before meshing.
\item Compile geometry and mesh when a solve or mesh picture requires them.
      A FEM-only export needs neither a mesh nor a solution.
\item Solve explicitly or through a field picture/profile. Capture measurement
      contours without turning them into physical interfaces.
\item Run the queued exports at picture end. If solving occurs later, use a
      named |\femmexport| command after that solve.
\end{enumerate}

Changing a source path, geometry tolerance, material or boundary requires a new
model and solve. The writer rejects a stale solution whose input signature has
changed. Cache format~5 validates geometry, physical settings and carved-domain
coverage; imported solutions restored from a matching cache can be exported.
Drawing scales, styles and profile sampling do not invalidate the physical cache.
Neither a FEM nor an ANS file reconstructs the original TeX program.

\subsection{Two U-shaped devices, in both directions}

The coil example uses the familiar node component. We add a depth and select the
FEMM material law before declaring the component, then request both exports at
picture end. The final drawing imports the produced FEM file and solves again.
The two coil regions are exported with their effective current densities. Their
original separate turn count and current cannot be recovered from an $NI$ value.
\luafemmexample{interchange-coil}

The permanent-magnet example follows the same sequence. The alnico region carries
a constant magnetisation direction; the surrounding iron remains nonlinear.
The field profiles belong to the original model even after the final imported
picture creates another model.
\luafemmexample{interchange-alnico}

\subsection{Supported features and explicit limits}

The reader accepts LF/CRLF, case-insensitive section keywords, quoted names,
optional label expressions as passive data, and all six FEMM length units.
Names remain opaque byte strings in the Lua codec: no Windows code-page to UTF-8
conversion is attempted. Numerical geometry rendering does not typeset those
names. The writer rejects quotes/newlines that cannot be represented safely in a
quoted scalar. Malformed input is diagnosed with filename and line number.

The solver adapter supports closed planar graphs, circular arcs up to $180^\circ$
per record, nested and disconnected regions, default labels, isotropic linear
and B--H materials, constant coercive direction, prescribed real current density,
series circuits, tagged prescribed-potential and constant real mixed conditions,
and natural boundaries. No Mesh seeds exclude their faces entirely. Each
unanchored connected component receives its own gauge and current-balance check.

The adapter rejects AC, axisymmetry, previous-solution modes, anisotropy,
unsupported lamination/fill physics, parallel circuits, point properties,
functional magnetisation, periodic/antiperiodic/sliding-gap boundaries, negative
Robin coefficients, and mixed conditions on internal interfaces. Open or isolated
constraints that do not bound valid faces also receive diagnostics. Default labels
with circuits, conflicting labels and labels on interfaces are rejected.
Unspecified FEM boundaries retain natural conditions; they do not inherit the
native rectangle's default prescribed potential.

Exact export rejects arbitrary source/boundary callbacks and affine mixed data
that a constant FEMM boundary property cannot represent. Segment groups are
retained for drawing. The computational export uses effective material records
per face and current density rather than claiming to recover original circuits
or every preprocessor editing attribute. For an unmodified passive document,
|luafemm-fem.serialize| retains its original property and arc records separately.

\subsection{Lua interface and independent validation}

The principal Lua entry points are:
\begin{itemize}
\item |luafemm.import_fem(file, options)| creates an unsolved model;
\item |luafemm.export_fem(model, file)| writes its computational geometry;
\item |luafemm.export_ans(model, file)| writes its converged solution.
\end{itemize}
The passive codec provides |read|, |parse|, |serialize| and |model|; parsing alone
does not authorize unsupported physics for solving. Imported expressions are
never passed to |load|, |dofile| or a TeX input command. The existing ownership
rule applies: imported documents and material tables are copied into the model.

Original small fixtures test units, indexing, malformed input, arc subdivision,
exclusions, labels, series-current normalization, cubic derivatives and cache
restoration. A separate optional C++ client loads emitted ANS files from disk
through xfemm's FPProc. It compares $A$, unsmoothed $B$ and $H$ with Lua samples.
Independent xfemm solves of the exported U devices are compared at two Lua mesh
resolutions; different meshers are not expected to produce identical triangles.
Run |python3 scripts/interop.py --build --xfemm ../xfemm| to reproduce these checks.
The audited revision and local compiler are recorded with the generated reports.
Windows FEMM GUI acceptance remains unverified on the macOS development host.
